% ====================================================================
% Apéndice B: Listados de código clave
% ====================================================================
\chapter{Listados de código clave}
\label{ap:listings}

Este apéndice reproduce los fragmentos clave del código fuente de 
\texttt{Mpemba-X v2}. El código completo está disponible en el archivo 
\texttt{mpemba-x-v2.tar.gz}.

\section{El parser INI (extracto de \texttt{core/config\_parser.hpp})}

\begin{lstlisting}[language=C++]
class Config {
public:
    void load(const std::string& path) {
        std::ifstream ifs(path);
        if (!ifs) throw std::runtime_error("Config: cannot open " + path);
        std::string line, section = "default";
        while (std::getline(ifs, line)) {
            line = trim(strip_comment(line));
            if (line.empty()) continue;
            if (line.front() == '[' && line.back() == ']') {
                section = trim(line.substr(1, line.size() - 2));
                continue;
            }
            auto eq = line.find('=');
            std::string key = trim(line.substr(0, eq));
            std::string value = trim(line.substr(eq + 1));
            data_[section][key] = value;
        }
    }
    
    double require_double(const std::string& s, const std::string& k) const {
        return std::stod(require_string(s, k));
    }
    
    std::vector<double> get_double_list(const std::string& s, 
                                         const std::string& k) const {
        std::vector<double> out;
        std::stringstream ss(get_string(s, k));
        std::string tok;
        while (std::getline(ss, tok, ','))
            if (!trim(tok).empty()) out.push_back(std::stod(tok));
        return out;
    }
private:
    std::map<std::string, std::map<std::string, std::string>> data_;
};
\end{lstlisting}

\section{Núcleo de la ecuación maestra markoviana}

\begin{lstlisting}[language=C++]
void rk4_step_dense(const DenseMarkov& M, std::vector<double>& p, double dt) {
    int n = M.n;
    std::vector<double> k1(n), k2(n), k3(n), k4(n), tmp(n);
    dense_matvec(M, p.data(), k1.data());
    
    #pragma omp parallel for schedule(static)
    for (int i = 0; i < n; ++i) tmp[i] = p[i] + 0.5 * dt * k1[i];
    dense_matvec(M, tmp.data(), k2.data());
    
    #pragma omp parallel for schedule(static)
    for (int i = 0; i < n; ++i) tmp[i] = p[i] + 0.5 * dt * k2[i];
    dense_matvec(M, tmp.data(), k3.data());
    
    #pragma omp parallel for schedule(static)
    for (int i = 0; i < n; ++i) tmp[i] = p[i] + dt * k3[i];
    dense_matvec(M, tmp.data(), k4.data());
    
    #pragma omp parallel for schedule(static)
    for (int i = 0; i < n; ++i)
        p[i] += dt * (k1[i] + 2.0*k2[i] + 2.0*k3[i] + k4[i]) / 6.0;
    
    // Renormalize to handle numerical drift
    double s = 0.0;
    for (int i = 0; i < n; ++i) { if (p[i] < 0) p[i] = 0.0; s += p[i]; }
    if (s > 0) for (int i = 0; i < n; ++i) p[i] /= s;
}
\end{lstlisting}

\section{Integrador Langevin sobreamortiguado}

\begin{lstlisting}[language=C++]
// Euler-Maruyama for overdamped Langevin
// dx = F(x) dt + sqrt(2 T dt) * eta
double sqrt_2T_dt = std::sqrt(2.0 * T_bath * dt);

for (int step = 0; step < n_steps; ++step) {
    #pragma omp parallel
    {
        int tid = omp_get_thread_num();
        PhiloxRNG& rng = rngs[tid];
        #pragma omp for schedule(static)
        for (int i = 0; i < N_traj; ++i) {
            double F  = potential.force(X[i]);
            double eta = rng.normal();
            X[i] += F * dt + sqrt_2T_dt * eta;
        }
    }
    // Histogram computation every n_steps_per_log steps...
}
\end{lstlisting}

\section{Esquema de volúmenes finitos para ecuación de Fourier}

\begin{lstlisting}[language=C++]
// Micro-step of the water Fourier model (Zhang 2014)
void micro_step(WaterColumn& w, double dt) {
    int N = w.N;
    double dx = w.dx;
    double tau = tau_HB(w.theta_initial);
    std::vector<double> theta_new(N), d_OH_new(N);
    
    #pragma omp parallel for schedule(static)
    for (int i = 0; i < N; ++i) {
        double q_in_left, q_out_right;
        
        // Robin BC at x=0 (drain)
        if (i == 0) {
            q_in_left = w.h_drain * (w.theta_bath - w.theta[i]);
        } else {
            double k_i = local_kappa(w, i), k_im = local_kappa(w, i-1);
            double kf = 2.0 * k_i * k_im / (k_i + k_im + 1e-30);
            q_in_left = -kf * (w.theta[i] - w.theta[i-1]) / dx;
        }
        
        // Adiabatic BC at x=L
        if (i == N - 1) {
            q_out_right = 0.0;
        } else {
            double k_i = local_kappa(w, i), k_ip = local_kappa(w, i+1);
            double kf = 2.0 * k_i * k_ip / (k_i + k_ip + 1e-30);
            q_out_right = -kf * (w.theta[i+1] - w.theta[i]) / dx;
        }
        
        // H-bond memory term
        double dOH_eq = d_OH_equilibrium(w.theta[i]);
        double ddOH = -(w.d_OH[i] - dOH_eq) / tau;
        double q_mem = w.C_mem * ddOH;
        
        double rhoCp = local_rho(w, i) * Cp_water(w.theta[i]);
        theta_new[i] = w.theta[i] + dt * (
            (q_in_left - q_out_right) / (dx * rhoCp) + q_mem / rhoCp);
        d_OH_new[i] = w.d_OH[i] + dt * ddOH;
    }
    w.theta.swap(theta_new);
    w.d_OH.swap(d_OH_new);
}
\end{lstlisting}

\section{Ejecución típica}

\begin{lstlisting}[language=bash, basicstyle=\ttfamily\footnotesize]
# Compilacion completa
mkdir build && cd build
cmake .. -DMPI_CXX_COMPILER=/usr/bin/mpicxx.openmpi
make -j4

# Ejecucion (uno por modulo)
OMP_NUM_THREADS=4 mpirun --allow-run-as-root --oversubscribe -np 2 \
    ./build/mpemba_water_fourier configs/water_demo.ini

# Visualizacion
python3 analysis/plot_water_fourier.py results/water_demo
\end{lstlisting}
